\begin{pastexam}[label={px:ncu-115-dsa-14}]{中央資工所 115 學年度 第 14 題}
（複選；每題 5 分，全對才給分，答錯倒扣 1 分）

Given list [9, 7, 15, 6, 11, 4, 22, 119]. Which of the following statements is (are) true?
\begin{enumerate}[label=\Alph*.]
\item The resulting sequence of the first phase of Quick Sort when the first element, i.e., 9, is chosen as pivot is 6, 7, 4, 9, 11, 15, 22, 119.
\item The resulting sequence of the first phase of Quick Sort when the second element, i.e., 7, is chosen as pivot is 4, 6, 7, 15, 11, 9, 22, 119.
\item At the end of the second pass of LSD Radix sort, the first element of the resulting chain is 11.
\item At the end of the third pass of LSD Radix sort, the 7-th element of the resulting chain is 22.
\end{enumerate}
\end{pastexam}

\begin{pastsoln}
答案：A、D。

題目沒有指定分割的方法。 「第一個當基準」最常見的是下面這種國內課本常用的做法：和資結書
虛擬碼~14.10 的 Hoare 分割一樣從兩頭往中間找，$i$ 從左往右找 $\ge x$ 的，$j$ 從右往左找 $\le x$
的，兩者交換；差別是基準不參與交換，等 $i$、$j$ 交錯時，才把基準和 $A[j]$ 交換，放到中間。 基準不在第一格時（B），先把它和第一格
交換，再用同一種分割。 基數排序用十進位，每一趟依一位數字分到 $10$ 個桶，桶內保持原來的順序
（資結書第~15.2.2~節）。

\begin{pseudo}{本題採用的分割}
\Function{Partition-First}{$A, p, r$}
    \State $x \gets A[p]$；$i \gets p + 1$；$j \gets r$ \Comment{基準是第一個}
    \While{\TRUE}
        \While{$i \le r$ \AND $A[i] < x$}
            \State $i \gets i + 1$
        \EndWhile
        \While{$A[j] > x$}
            \State $j \gets j - 1$
        \EndWhile
        \If{$i \ge j$}
            \State 交換 $A[p]$ 和 $A[j]$；\Return $j$
        \EndIf
        \State 交換 $A[i]$ 和 $A[j]$；$i \gets i + 1$；$j \gets j - 1$
    \EndWhile
\EndFunction
\end{pseudo}

\textbf{A 對}。 位置從 $1$ 編號，$x = 9$：
\begin{itemize}
\item $i$ 從第 $2$ 格往右：$7 < 9$，停在第 $3$ 格的 $15$。 $j$ 從第 $8$ 格往左：$119$、$22$ 都 $> 9$，
停在第 $6$ 格的 $4$。 交換，得 $\langle 9, 7, 4, 6, 11, 15, 22, 119\rangle$，$i = 4$、$j = 5$。
\item $i$：$6 < 9$，停在第 $5$ 格的 $11$。 $j$：第 $5$ 格 $11 > 9$，停在第 $4$ 格的 $6$。
\item $i = 5 \ge j = 4$，把基準和第 $4$ 格交換：$\langle 6, 7, 4, 9, 11, 15, 22, 119\rangle$。
\end{itemize}
和 A 完全相同。

\textbf{B 錯}。 先把 $7$ 換到第一格：$\langle 7, 9, 15, 6, 11, 4, 22, 119\rangle$，$x = 7$。
\begin{itemize}
\item $i$ 停在第 $2$ 格的 $9$；$j$ 跳過 $119$、$22$，停在第 $6$ 格的 $4$。 交換：
$\langle 7, 4, 15, 6, 11, 9, 22, 119\rangle$。
\item $i$ 停在第 $3$ 格的 $15$；$j$ 跳過 $11$，停在第 $4$ 格的 $6$。 交換：
$\langle 7, 4, 6, 15, 11, 9, 22, 119\rangle$。
\item $i$ 停在第 $4$ 格的 $15$；$j$ 停在第 $3$ 格的 $6$。 $i \ge j$，基準和第 $3$ 格交換，得到
$\langle 6, 4, 7, 15, 11, 9, 22, 119\rangle$。
\end{itemize}
右邊 $15, 11, 9, 22, 119$ 和 B 一樣，但左邊是 $6, 4$，不是 $4, 6$。 改用其他常見的分割也得不到 B 的
序列：基準換到最後再用 Lomuto 分割（資結書虛擬碼~14.8），得 $\langle 6, 4, 7, 9, 11, 119, 22, 15\rangle$；
基準留在第一格、由左往右只掃一次的 Lomuto 寫法，得 $\langle 4, 6, 7, 9, 11, 15, 22, 119\rangle$。
所以 B 錯。

\textbf{LSD 基數排序}。 最大的數 $119$ 有三位，做三趟（個位、十位、百位），每一趟依序收集
$0$ 號到 $9$ 號桶：
\begin{center}
\small
\begin{tabular}{lll}
\toprule
趟 & 非空的桶 & 收集後的串列 \\
\midrule
1（個位）& 1:\,11；2:\,22；4:\,4；5:\,15；6:\,6；7:\,7；9:\,9, 119 & 11, 22, 4, 15, 6, 7, 9, 119 \\
2（十位）& 0:\,4, 6, 7, 9；1:\,11, 15, 119；2:\,22 & 4, 6, 7, 9, 11, 15, 119, 22 \\
3（百位）& 0:\,4, 6, 7, 9, 11, 15, 22；1:\,119 & 4, 6, 7, 9, 11, 15, 22, 119 \\
\bottomrule
\end{tabular}
\end{center}
第一趟 $9$ 號桶裡 $9$ 在 $119$ 前面，因為原串列裡 $9$ 先出現；第二趟 $119$ 的十位數是 $1$，和 $11$、
$15$ 同桶，排在它們後面。

\textbf{C 錯}。 第二趟結束時串列的第一個是 $4$，不是 $11$。

\textbf{D 對}。 第三趟結束時已經排好，第 $7$ 個是 $22$。
\end{pastsoln}
